\manualchapter{Overview and first sound}{sec:quick-start}
RackNES runs an NES game inside VCV Rack 2. The game supplies the music,
sound effects, and visuals; your patch controls the players, changes the
emulation speed, and returns to saved moments. Five separate sound outputs
and a mix output connect the game to the rest of your rack.

The module starts without a game. Supply an iNES ROM file and use the panel
buttons or gates as the two NES controllers. A game's response to those
buttons depends on its menus and current play state.

\subsection{Start here}
\begin{enumerate}
\item Add RackNES and connect \textbf{MIX} to your audio output through a
mixer or attenuator. Begin with a low listening level.
\item Leave all five channel levels and Mix Volume at their default
\textbf{100\%}. Leave the large clock knob at its default
\textbf{1.789773 MHz} and the red clock attenuverter at \textbf{0\%}.
\item Right-click RackNES, choose \textbf{Load ROM}, and select a
\texttt{.nes} or \texttt{.NES} file. You can also drop one ROM file onto the
module. The supported cartridge mappers are listed on page~\pageref{sec:roms}.
\item Use Player 1 \textbf{Start}, \textbf{Select}, the directions, and
\textbf{A/B} to reach a part of the game that makes sound. Hold a panel
button for as long as the game requires.
\item Press \textbf{SAVE} at an interesting moment. Let the game continue,
then press \textbf{LOAD} to revisit it. There is one snapshot slot.
\end{enumerate}

\subsection{Three useful first experiments}
\begin{description}
\item[Separate the percussion.] Patch \textbf{NOI} to a second mixer input.
Noise disappears from MIX and can now have its own effects chain.
\item[Automate a controller.] Send a repeating gate to Player 1 A.
A high gate holds the button; a low gate releases it.
\item[Repeat a game moment.] Send a slow trigger to LOAD after making a
snapshot. Each trigger restores the stored emulation state.
\end{description}

\subsection{What to expect}
RackNES is an emulator under voltage control. Its clock input changes the
speed of execution; it does not receive tempo pulses for synchronization.
SAVE and LOAD store machine state, not an audio clip. Hang freezes execution
and holds the outputs. The following sections explain these distinctions
and the limitations of the current implementation.
